\section{1D Convolutional Neural Networks: Theoretical Background}
This section provides a comprehensive, ground-up theoretical foundation for 1D Convolutional Neural Networks (1D-CNNs). It is designed to provide sufficient mathematical rigor so that a reader could implement a 1D-CNN from scratch in C or NumPy, without relying on high-level libraries like PyTorch.

\subsection{The 1D Convolution Operation (Forward Pass)}
Standard fully-connected networks are ill-suited for time-series signal processing. They treat the input as an unordered set of independent variables, destroying temporal topology, and they lack \textbf{translation equivariance} (the ability to recognize a pattern regardless of where it occurs in time).

A 1D-CNN solves this by using a "sliding window" approach. Instead of learning a massive weight matrix for the entire signal, the network learns small templates called \textbf{filters} or \textbf{kernels}.

\subsubsection{Single Channel Convolution}
Let $\mathbf{x}$ be a 1D input signal of length $T$, and $\mathbf{w}$ be a learnable kernel of size $K$. The convolution operation slides the kernel $\mathbf{w}$ across the signal $\mathbf{x}$. At each time step $t$, it computes the dot product between the kernel and the local patch of the signal:
\begin{equation}
    y[t] = \sum_{i=0}^{K-1} w[i] \cdot x[t + i] + b
\end{equation}
where $b$ is a learnable bias scalar. Mathematically, this dot product acts as a \textbf{pattern matcher}. If the signal patch $x[t \dots t+K-1]$ is perfectly aligned with the weights in the kernel $w$, the output $y[t]$ produces a massive positive activation. The resulting sequence $\mathbf{y}$ is called a \textbf{feature map}.

\subsubsection{Multi-Channel Convolution}
A single filter can only detect one specific shape (e.g., a rising edge). To detect complex anomalies, a layer must learn dozens of filters simultaneously. Furthermore, deeper layers in the network do not operate on the raw 1D signal; they operate on the multiple feature maps generated by the previous layer.

Let $C_{in}$ be the number of input channels and $C_{out}$ be the number of desired output channels (filters). The kernel is now a 3D tensor $\mathbf{W}$ of shape $(C_{out}, C_{in}, K)$. The output feature map for a specific output channel $c_{out}$ at time $t$ is the sum of convolutions across all input channels $c_{in}$:
\begin{equation}
    y_{c_{out}}[t] = \sum_{c_{in}=0}^{C_{in}-1} \sum_{i=0}^{K-1} W[c_{out}, c_{in}, i] \cdot x_{c_{in}}[t + i] + b_{c_{out}}
\end{equation}

This operation preserves \textbf{local connectivity} (neurons only look at a small temporal patch $K$) and enforces \textbf{weight sharing} (the exact same tensor $\mathbf{W}$ is applied at every time step $t$), which drastically prevents overfitting and reduces computational cost.

\subsection{Non-linear Activations}
Without a non-linear activation function, stacking multiple convolutional layers would mathematically collapse into a single, linear convolution. After computing the feature map $\mathbf{y}$, we apply an activation function $\sigma$ element-wise.

\subsubsection{ReLU (Rectified Linear Unit)}
The most widely used activation in CNN hidden layers is ReLU. It passes positive activations (where the filter "found" its target pattern) and zeroes out negative ones:
\begin{equation}
    a[t] = \max(0, y[t])
\end{equation}
Its derivative is trivially $1$ if $y[t] > 0$, and $0$ otherwise. This simplicity prevents the vanishing gradient problem during deep backpropagation.

\subsubsection{Softmax}
At the very end of the network, after flattening the final feature maps into a vector of logits $\mathbf{z}$ (one for each class), we apply the Softmax function to convert them into a valid probability distribution:
\begin{equation}
    \hat{y}_c = \frac{e^{z_c}}{\sum_k e^{z_k}}
\end{equation}

\subsection{Pooling Layers (Downsampling)}
To hierarchically build larger receptive fields and reduce computational cost, CNNs use \textbf{Pooling layers} between convolutions. 

\textbf{MaxPooling} slides a small window of size $P$ (usually 2) across the feature map with a stride of $P$. Instead of a dot product, it simply outputs the maximum value in that window:
\begin{equation}
    p[t] = \max_{i \in [P \cdot t, P \cdot t + P - 1]} a[i]
\end{equation}

Intuitively, MaxPooling asks: \textit{"Did this specific feature appear anywhere within this small region?"} If yes, it keeps the high activation and discards the exact spatial coordinate. This halves the length of the signal ($T \rightarrow T/2$) and provides \textbf{local translation invariance} against small signal jitters. 

Crucially for implementation, during the forward pass, a MaxPooling layer must cache the \textbf{argmax index} (the exact position $i$ that contained the maximum value) because gradients will only be routed back through that specific winning "pixel" during backpropagation.

\subsection{The Loss Function}
To train the CNN, we quantify its prediction error using \textbf{Cross-Entropy Loss}. For a single window, comparing the true one-hot label vector $\mathbf{y}$ against the predicted probabilities $\mathbf{\hat{y}}$:
\begin{equation}
    \mathcal{L} = -\sum_k y_k \log(\hat{y}_k)
\end{equation}

\subsubsection*{Proof: Derivative of Softmax with Cross-Entropy}
A fundamental result in machine learning is how elegantly the gradient simplifies at the output layer. We want to find the error signal $\delta_j = \frac{\partial \mathcal{L}}{\partial z_j}$ where $z_j$ is the pre-softmax logit. By the multivariate chain rule:
\begin{equation}
    \frac{\partial \mathcal{L}}{\partial z_j} = \sum_k \frac{\partial \mathcal{L}}{\partial \hat{y}_k} \frac{\partial \hat{y}_k}{\partial z_j}
\end{equation}

1. Derivative of Loss: $\frac{\partial \mathcal{L}}{\partial \hat{y}_k} = -\frac{y_k}{\hat{y}_k}$
2. Derivative of Softmax (quotient rule): 
   - If $k = j$: $\hat{y}_j(1 - \hat{y}_j)$
   - If $k \neq j$: $-\hat{y}_k \hat{y}_j$

Substituting these into the sum:
\begin{equation*}
    \frac{\partial \mathcal{L}}{\partial z_j} = \left( -\frac{y_j}{\hat{y}_j} \right) \big( \hat{y}_j(1 - \hat{y}_j) \big) + \sum_{k \neq j} \left( -\frac{y_k}{\hat{y}_k} \right) \big( -\hat{y}_k \hat{y}_j \big)
\end{equation*}
\begin{equation*}
    = -y_j(1 - \hat{y}_j) + \sum_{k \neq j} y_k \hat{y}_j = -y_j + y_j \hat{y}_j + \hat{y}_j \sum_{k \neq j} y_k
\end{equation*}

Since the true labels $y$ must sum to 1, $\sum_{k \neq j} y_k = 1 - y_j$. Substituting this:
\begin{equation*}
    = -y_j + y_j \hat{y}_j + \hat{y}_j (1 - y_j) = -y_j + y_j \hat{y}_j + \hat{y}_j - y_j \hat{y}_j
\end{equation*}

The complex terms cancel out, leaving the beautiful final result for the output layer error:
\begin{equation}
    \delta_j^{\text{out}} = \frac{\partial \mathcal{L}}{\partial z_j} = \hat{y}_j - y_j
\end{equation}

\subsection{Backpropagation in 1D-CNNs (Backward Pass)}
Once the output error $\delta^{\text{out}}$ is calculated, the \textbf{Backpropagation algorithm} uses the chain rule to pass this error backward through the network, layer by layer, to compute the gradient of the loss with respect to every single filter weight.

\subsubsection{Backpropagating through MaxPooling}
Because MaxPooling only routed the maximum value forward, it acts as a gate. The error $\delta$ from the subsequent layer is routed entirely to the specific index that "won" the max pooling (which we cached during the forward pass). All other indices in the pooling window receive a gradient of 0, because changing them slightly would not have affected the forward output.

\subsubsection{Backpropagating through 1D Convolution}
This is the mathematical core of training a CNN. Let the error signal received from the next layer be $\delta[t] = \frac{\partial \mathcal{L}}{\partial y[t]}$. We must compute two sets of gradients for our single-channel convolution equation (Eq. 1).

\textbf{1. Gradient with respect to the filter weights ($w[i]$):}
A single weight $w[i]$ in the filter was multiplied by the input $x[t+i]$ at \textit{every} time step $t$ as the filter slid across the signal. By the chain rule, we sum its contribution over all time steps:
\begin{equation}
    \frac{\partial \mathcal{L}}{\partial w[i]} = \sum_t \frac{\partial \mathcal{L}}{\partial y[t]} \frac{\partial y[t]}{\partial w[i]} = \sum_t \delta[t] \cdot x[t+i]
\end{equation}
This reveals a profound mathematical symmetry: to find how to update the filter weights, we perform a \textbf{cross-correlation} between the forward input signal $x$ and the backward error signal $\delta$. (For multi-channel convolutions, this cross-correlation is summed across all batch items and input dimensions).

\textbf{2. Gradient with respect to the input ($x[t]$):}
To pass the error further backward to earlier layers, we need $\frac{\partial \mathcal{L}}{\partial x[t]}$. A single input point $x[t]$ affects multiple outputs $y$ as the kernel slides over it. Specifically, it is multiplied by $w[0]$ to create $y[t]$, $w[1]$ to create $y[t-1]$, and so on up to $w[K-1]$ to create $y[t-(K-1)]$.

Summing these contributions via the chain rule:
\begin{equation}
    \frac{\partial \mathcal{L}}{\partial x[t]} = \sum_{i=0}^{K-1} \frac{\partial \mathcal{L}}{\partial y[t-i]} \frac{\partial y[t-i]}{\partial x[t]} = \sum_{i=0}^{K-1} \delta[t-i] \cdot w[i]
\end{equation}
This is an elegant symmetry: to pass the error backward, we perform a \textbf{full convolution} of the error signal $\delta$ with the \textbf{flipped} kernel $w$. This allows deep learning frameworks to reuse highly optimized forward convolution algorithms for the backward pass.

\subsection{Optimization and Regularization}
\subsubsection{Adam Optimizer}
With the gradient of the weights $\nabla_\mathbf{w} \mathcal{L}$ computed, we update the filters to minimize the loss. Instead of pure Stochastic Gradient Descent ($\mathbf{w} \leftarrow \mathbf{w} - \eta \nabla_\mathbf{w} \mathcal{L}$), we use the \textbf{Adam} optimizer. Adam maintains exponential moving averages of both the gradients (momentum) and the squared gradients (scaling). This provides an adaptive, per-parameter learning rate that navigates flat plateaus and steep ravines in the loss landscape much faster than standard SGD.

\subsubsection{Regularization: Dropout and BatchNorm}
To prevent the CNN filters from memorizing the specific noise patterns of the training data (overfitting), two mechanisms are used:
\begin{itemize}
    \item \textbf{Dropout:} Randomly zeroes out a percentage of the activations in the fully connected classification head during training. This prevents the network from over-relying on any single feature detector.
    \item \textbf{Batch Normalization:} Inserted immediately after convolutional layers, it normalizes the output feature maps to have zero mean and unit variance across the mini-batch. This stabilizes the distribution of inputs to the next layer, allowing for higher learning rates and significantly faster convergence.
\end{itemize}

\newpage
